\section{Introdução}
\label{sec:introducao}

\subsection{Contexto e motivação}

O \emph{Vehicle Routing Problem with Time Windows} (VRPTW) combina decisões de roteamento com restrições temporais: além de visitar todos os clientes, cada atendimento deve ocorrer dentro de uma janela de tempo predefinida. Essa combinação torna o problema significativamente mais difícil do que variantes sem restrições de tempo, tanto do ponto de vista de modelagem quanto de resolução computacional \cite{toth2002vehicle,braysy2005vehicle}.

Na prática, o VRPTW aparece em distribuição urbana, abastecimento de lojas, entregas com horário agendado e operações de coleta. Nesses contextos, soluções de boa qualidade precisam equilibrar custo de deslocamento, uso de frota e viabilidade operacional. Em instâncias de porte realista, métodos exatos tendem a ter custo elevado, o que justifica o uso de heurísticas e meta-heurísticas bem estruturadas.

\subsection{Base acadêmica adotada}

Este trabalho segue a linha clássica da literatura: construção inicial por Solomon I1 \cite{solomon1987algorithms} e refinamento por busca local/meta-heurística. Três métodos de melhoria foram implementados e avaliados:
\begin{itemize}
  \item VND (\emph{Variable Neighborhood Descent}) \cite{hansen2001variable};
  \item GRASP nas variantes fixa e reativa \cite{feo1995greedy,resende2003greedy,prais2000reactive};
  \item Busca Tabu clássica com memória de curto prazo e aspiração \cite{glover1989tabu,glover1990tabu}.
\end{itemize}

A proposta metodológica central foi manter a comparação justa: o construtivo base é o mesmo para todos os algoritmos, com função objetivo única (custo total), critérios de factibilidade idênticos e protocolo de execução reprodutível.

\subsection{Objetivo do trabalho}

O objetivo é duplo:
\begin{enumerate}
  \item garantir alinhamento entre definição acadêmica e lógica implementada no código;
  \item comparar o comportamento de VND, GRASP e Tabu quando todos partem do mesmo ponto de entrada construtivo.
\end{enumerate}

Em termos práticos, isso exige explicitar não apenas os pseudocódigos, mas também as estruturas de dados, as regras de aceitação de movimentos, os critérios de parada e os parâmetros passíveis de calibração.

\subsection{Contribuições deste texto}

As principais contribuições documentadas são:
\begin{itemize}
  \item padronização formal da notação matemática e das variáveis usadas nos algoritmos;
  \item descrição detalhada do fluxo de execução do solver e das estruturas de dados centrais;
  \item especificação do Solomon I1 na configuração clássica orientada à distância;
  \item análise de VND com seleção de 4 vizinhanças a partir de evidência empírica;
  \item descrição unificada de GRASP (fixo e reativo) e Busca Tabu clássica;
  \item protocolo de calibração por \emph{irace} e discussão de reprodutibilidade experimental.
\end{itemize}

\subsection{Organização do artigo}

A Seção~\ref{sec:notacao_fluxo} apresenta notação, estruturas de dados e fluxo geral do solver. A Seção~\ref{sec:vrptw_i1} descreve o modelo do problema e o construtivo Solomon I1. As Seções~\ref{sec:vnd}, \ref{sec:grasp} e \ref{sec:tabu} detalham os métodos de melhoria. A Seção~\ref{sec:protocolo} consolida protocolo experimental, calibração e resultados preliminares. Por fim, a Seção~\ref{sec:conclusao} apresenta síntese e próximos passos.
